% =====================================================
% Abstract
% Target length: 150-200 words (per AECD template requirement).
% Quantitative results are intentionally left as placeholders --
% see chapter08_experimental_results.tex and chapter06_multisim_simulation.tex
% for where the team must insert the final measured/simulated figures.
% =====================================================
\chapter*{Abstract}
\addcontentsline{toc}{chapter}{Abstract}

This report presents the design, simulation, and hardware implementation of an Automatic Gain Control (AGC) microphone preamplifier, developed for the course 25EC1302E -- Analog Electronic Circuit Design. Microphone signal levels vary widely with speaker distance and vocal intensity; a fixed-gain amplifier therefore either clips on loud input or gives poor signal-to-noise performance on quiet input. The proposed circuit addresses this by combining a low-noise preamplifier stage with a variable-gain feedback loop that senses the output envelope and adjusts amplifier gain so the output stays within a target level range. This report presents the operating principle, the design methodology and governing equations, the Multisim simulation plan, and the breadboard-based hardware verification procedure used to develop the circuit. The final mid-band gain, AGC threshold, compression ratio, attack and release times, and measured signal-to-noise performance are \placeholder{insert final measured and simulated quantitative results} pending completion of simulation and bench measurement by the project team. Key limitations of the chosen approach and low-voltage prototype safety considerations are also summarised.

\vspace{0.5cm}
\noindent\textbf{Keywords:} automatic gain control; microphone preamplifier; variable-gain amplifier; envelope detection; analog circuit design; Multisim; \placeholder{add 1--2 project-specific keywords, e.g. the gain-control IC/topology used}.
